\section{System Overview}
\label{sec:overview}

\subsection{Design principles}

The implementation follows six principles.

\begin{description}[leftmargin=1.4em,style=nextline]
  \item[Fail closed.] The decoder accepts only the documented subset. A branch
        target outside the function, an overlapping instruction, an unknown flag
        producer, an unsupported prefix, or a malformed encoding is a build
        error. The runtime rejects malformed payloads, invalid registers,
        out-of-range memory, invalid opcodes, and exhausted fuel with a
        deterministic status. There is no best-effort translation.
  \item[Zero shipped dependencies.] The compiler core, VM, and binary-format
        layer depend only on \code{core}. The protector CLI uses \code{std} at
        the filesystem, process, and error-reporting boundary. No third-party
        crate is linked into a protected artifact.
  \item[Memory safety by construction.] The safe tier uses
        \texttt{\#![forbid(unsafe\_code)]}. The VM has no heap allocation, no
        unchecked indexing in its public path, no division, no floating point,
        and no host-memory access. Bounds are checked at every decode and memory
        operation.
  \item[Per-build diversity.] A 256-bit seed drives opcode assignment, alias
        operand order, register-map selection, immediate endianness, manifest
        endianness, state identifiers, block order, rolling keys, entry shares,
        MBA shape, decoy placement, semantic-handler order, and code/data
        offsets. Two builds of the same input differ in every payload byte and
        in most of the VM segment.
  \item[Separated unsafe boundary.] Debugger detection, RDTSC/LFENCE, and
        native text checksums require OS interfaces or instructions that safe
        Rust cannot express. They are isolated in an opt-in
        \code{sirus-hardened} crate and generated only with \verb|--hardened|.
  \item[Evidence rather than assertion.] Every security claim is paired with a
        script that can be re-executed. The paper reports negative results and
        tool limitations alongside positive ones.
\end{description}

\subsection{Build pipeline}

Figure~\ref{fig:pipeline} shows the build pipeline. The input image is parsed and
validated; one or more selectors resolve functions; the selected ranges are
decoded and lifted; the resulting CFG is flattened and rewritten; the IR is
encoded into the randomized ISA; the manifest and bodies are encrypted; a
freestanding Rust VM is generated and compiled; the resulting blob is validated;
the blob is injected into a new executable segment; the original entries are
patched; and the original bodies are scrubbed.

\begin{figure}[htbp]
\centering
\begin{tikzpicture}[
  box/.style={draw, rounded corners=1.5pt, align=center, text width=30mm,
              minimum height=9mm, font=\small, fill=white},
  arrow/.style={-{Latex[length=2mm]}, thick, color=sirusblue},
  node distance=6mm and 7mm
]
\matrix (m) [matrix of nodes, nodes={box}, row sep=7mm, column sep=7mm] {
  Input ELF/PE & Parse and validate & Resolve selectors & Decode \\
  Lift to IR & Build CFG & Flatten & MBA, decoys, blinding \\
  Encode randomized ISA & Encrypt manifest/bodies & Synthesize no\_std VM & Compile and validate \\
  Inject RX segment & Patch entries & Scrub bodies & Protected output \\
};
\foreach \i [evaluate=\i as \j using int(\i+1)] in {1,2,3} {
  \draw[arrow] (m-1-\i) -- (m-1-\j);
  \draw[arrow] (m-2-\j) -- (m-2-\i);
  \draw[arrow] (m-3-\i) -- (m-3-\j);
  \draw[arrow] (m-4-\j) -- (m-4-\i);
}
\draw[arrow] (m-1-4) -- (m-2-4);
\draw[arrow] (m-2-1) -- (m-3-1);
\draw[arrow] (m-3-4) -- (m-4-4);
\end{tikzpicture}
\caption{Build pipeline. The flow snakes through four rows; each stage is
deterministic for a fixed input, seed, profile, and toolchain.}
\label{fig:pipeline}
\end{figure}

The target binary is never relocated. The ELF program-header table is moved to
an appended read-only page, one new read-only \code{PT_LOAD} maps it, and one
new read-execute \code{PT_LOAD} maps the VM. The original load segments are
untouched. The PE backend appends one read-execute section, updates the section
count and image size, recomputes the checksum, and patches RVA-relative entries.

\subsection{Runtime pipeline}

Figure~\ref{fig:runtime} shows the runtime. A host call reaches a per-function
ABI shim. The shim preserves the host calling convention, materializes register
and stack arguments, and calls the generated entry. The entry derives the
decode tables and payload key from the seed and the integrity value, scans the
encrypted manifest for the entry state, and enters the dispatcher. The dispatcher
decrypts one instruction at a time, resolves the alias, executes the semantic
operation, and transitions to the next state. A return writes the virtual
\code{RAX} value; an error writes zero and a status.

\begin{figure}[htbp]
\centering
\begin{tikzpicture}[
  box/.style={draw, rounded corners=1.5pt, align=center, text width=34mm,
              minimum height=9mm, font=\small, fill=white},
  smallbox/.style={draw, rounded corners=1.5pt, align=center, text width=28mm,
              minimum height=8mm, font=\small, fill=white},
  arrow/.style={-{Latex[length=2mm]}, thick, color=sirusblue},
  node distance=7mm and 12mm
]
\node[box] (call) {Host call to patched entry};
\node[box, below=of call] (shim) {ABI shim: preserve registers, copy arguments};
\node[box, below=of shim] (init) {Derive decode tables and payload key};
\node[box, below=of init] (manifest) {Scan encrypted manifest for entry state};
\node[box, below=of manifest] (decode) {Decrypt and decode one instruction};
\node[box, below=of decode] (exec) {Execute semantic operation};
\node[smallbox, right=of exec] (next) {Transition to next state};
\node[smallbox, above=of next] (ret) {Return virtual RAX};
\node[smallbox, below=of next] (fail) {BadPayload / Fuel / LoopBound};
\draw[arrow] (call) -- (shim);
\draw[arrow] (shim) -- (init);
\draw[arrow] (init) -- (manifest);
\draw[arrow] (manifest) -- (decode);
\draw[arrow] (decode) -- (exec);
\draw[arrow] (exec) -- (next);
\draw[arrow] (next) |- (decode);
\draw[arrow] (next) -- (ret);
\draw[arrow] (next) -- (fail);
\end{tikzpicture}
\caption{Runtime structure. The transition arrow is the dispatcher loop; the
manifest is scanned once per transition in the normal profile.}
\label{fig:runtime}
\end{figure}

\subsection{Crate architecture and trust boundaries}

The workspace separates concerns by trust level:

\begin{itemize}
  \item \code{sirus-core}: \code{no_std}, no allocation, no unsafe. Decoder,
        lifter, IR, MBA arena, ISA derivation, payload cipher, program builder.
  \item \code{sirus-vm}: \code{no_std}, no allocation, no unsafe. Reference
        engine used by tests and embedded verbatim into generated runtimes.
  \item \code{sirus-binfmt}: \code{no_std} plus \code{alloc}, no unsafe. ELF and
        PE parsers, editors, and injection geometry.
  \item \code{sirus-protector}: \code{std} boundary, safe Rust. Selector
        resolution, source generation, \code{rustc} invocation, blob validation,
        injection, patching, and scrubbing.
  \item \code{sirus-hardened}: opt-in unsafe. Linux and Windows hardening
        primitives; never linked into the safe default runtime.
\end{itemize}

The compiler and runtime share the \code{shared} modules verbatim: the ISA
constants, SipHash, ChaCha20, and decode-table derivation are compiled into both
the protector and the generated VM. This eliminates a class of encoder/decoder
divergence bugs and is checked by the build system, which tracks every embedded
source file and rebuilds the protector when one changes.

The external verification crates (\code{verification/dudect},
\code{verification/creusot}, and \code{fuzz}) are outside the shipped dependency
graph. They may use external crates because they are measurement and proof
tooling, not runtime components.

\subsection{Threat surface}

The protected binary has three relevant surfaces: the host image, the injected
VM segment, and the runtime entry path. The host image retains its symbol table
and dynamic metadata; Sirus does not strip the host. The VM segment contains the
engine, the encrypted payload, the encrypted decode table, and the per-build
constants. The entry path contains the patched jump and the ABI shim. An
attacker can read all three, trace them, and patch them. The design assumes this
and attempts to make the semantic recovery and patch construction expensive, not
impossible.
